\section*{Chapter \ref{ch:b2:biomolecules} --- Biomolecules: Amino Acids, Sugars, Lipids and Nucleotides}

\begin{solution}{exo:b2:biomolecules:1}
\begin{itemize}
  \item Between pH 3 and 9, the \omterm{def:b2:biomolecules:amino-acid}{zwitterion} \ce{H3N+-CH(CH3)-COO-}.
  \item At pH 1, the cation \ce{H3N+-CH(CH3)-COOH}.
  \item At pH 12, the anion \ce{H2N-CH(CH3)-COO-}.
\end{itemize}
\end{solution}

\begin{solution}{exo:b2:biomolecules:2}
Glycine: $(2.35 + 9.78)/2 \approx 6.1$. Alanine: $(2.35 + 9.87)/2 \approx 6.1$.
% ledger: pka:glycine1, pka:glycine2, pka:alanine1, pka:alanine2
\end{solution}

\begin{solution}{exo:b2:biomolecules:3}
\ce{H2N-CH(CH3)-CO-NH-CH2-COOH}: alanine provides the N-terminus, glycine the C-terminus. Gly--Ala,
\ce{H2N-CH2-CO-NH-CH(CH3)-COOH}, is a constitutional isomer: the free amino group is on glycine.
\end{solution}

\begin{solution}{exo:b2:biomolecules:4}
\ce{NH2} on the left: L. Priorities \ce{NH2} > \ce{COOH} > \ce{CH2OH} > \ce{H}: L-serine is
(\textit{S}).
\end{solution}

\begin{solution}{exo:b2:biomolecules:5}
Groups: N-terminal ammonium (about 9), lysine \omterm{def:b2:biomolecules:amino-acid}{side chain} (10.7), aspartate \omterm{def:b2:biomolecules:amino-acid}{side chain} (3.9),
C-terminal carboxylic group (about 2). pH 1: $+1 +1 + 0 + 0 = +2$. pH 7: $+1 +1 -1 -1 = 0$. pH 12:
$0 + 0 - 1 - 1 = -2$ (the lysine \omterm{def:b2:biomolecules:amino-acid}{side chain}, 1.3 units below pH 12, is mostly neutral).
\end{solution}

\begin{solution}{exo:b2:biomolecules:6}
Thionyl chloride and heat would also attack other groups and, above all, the very reactive \omterm{def:b2:acyl-substitution:derivative}{acyl chloride}
of an amino acid loses its configuration easily (the $\alpha$ hydrogen becomes acidic): the \omterm{def:b2:biomolecules:peptide}{peptide} would
be partly racemised. A \omterm{def:b2:biomolecules:peptide}{coupling agent} activates the acid at room temperature, under mild conditions.
\end{solution}

\begin{solution}{exo:b2:biomolecules:7}
O5 bonded to C2: a ring of C2, C3, C4, C5 and O, five atoms (a furanose). In \omterm{def:b2:biomolecules:anomer}{Haworth projection}, ring O
at the back, C2 at the right with \ce{CH2OH} (C1) down and \ce{OH} up ($\beta$); C3 \ce{OH} up (left in
Fischer), C4 \ce{OH} down, C5 carrying \ce{CH2OH} (C6) up.
\end{solution}

\begin{solution}{exo:b2:biomolecules:8}
$x_\alpha = (80.2 - 52.8)/(150.7 - 52.8) = 27.4/97.9 \approx 0.28$: 28~\% $\alpha$, 72~\% $\beta$.
\end{solution}

\begin{solution}{exo:b2:biomolecules:9}
Maltose keeps a free hemiacetal on its second glucose, which opens to an aldehyde: reducing. In sucrose
both \omterm{def:b2:biomolecules:anomer}{anomeric carbons} are engaged in the \omterm{def:b2:biomolecules:glycoside}{glycosidic bond}: no hemiacetal, not reducing. Dilute acid
hydrolyses the \omterm{def:b2:biomolecules:glycoside}{glycosidic bond}, an acetal: glucose and fructose are released, and both reduce.
\end{solution}

\begin{solution}{exo:b2:biomolecules:10}
Triolein, \qty{885.45}{g/mol}, consumes three \ce{KOH} (\qty{56.105}{g/mol}): $3 \times 56.105/885.45 =
0.190$~g per g, a \omterm{def:b2:acyl-substitution:saponification}{saponification} value of 190. Shorter chains mean a smaller molar mass for the same
three ester groups: more moles of fat per gram, more KOH.
% ledger: aw:C, aw:H, aw:O, aw:K
\end{solution}

\begin{solution}{exo:b2:biomolecules:11}
Protect the amine of glycine as a \textit{tert}-butoxycarbonyl carbamate, and the acid of alanine as a
methyl ester. Couple with a carbodiimide. Remove the carbamate with acid and the ester with dilute base
(orthogonal conditions, so either end can be freed alone): Gly--Ala.
\end{solution}

\begin{solution}{exo:b2:biomolecules:12}
5$'$-ACGCAT-3$'$ (read antiparallel: A pairs with T, G with C). Pairs: A--T, T--A, G--C, C--G, G--C, T--A:
$3 \times 2 + 3 \times 3 = 15$ hydrogen bonds.
\end{solution}

\begin{solution}{pb:b2:biomolecules:1}
\textbf{1.} L-Aspartic acid, L-phenylalanine and methanol.
\textbf{2.} In both, \ce{COOH} at the top, \omterm{def:b2:biomolecules:amino-acid}{side chain} at the bottom (\ce{CH2COOH}, \ce{CH2C6H5}), \ce{NH2}
on the left.
\textbf{3.} (\textit{S}) for both.
\textbf{4.} Two stereocentres: four stereoisomers.
\textbf{5.} \ce{H2N-CH(CH2COOH)-CO-NH-CH(CH2C6H5)-COOCH3}: the \omterm{def:b2:biomolecules:peptide}{peptide bond} joins the two
residues; the ester is the \ce{COOCH3} at the C-terminus.
\textbf{6.} Taste receptors are chiral: they bind one stereoisomer and not its mirror image or
diastereoisomers, as the opening said.
\textbf{7.} The N-terminal \ce{NH3+} and the side-chain \ce{COOH} of aspartic acid; the C-terminal acid
is a methyl ester, with no acidic hydrogen.
\textbf{8.} At pH 2: \ce{NH3+} ($+1$), side-chain \ce{COOH} (0): $+1$.
\textbf{9.} At pH 7: $+1 - 1 = 0$.
\textbf{10.} At pH 12: $0 - 1 = -1$.
\textbf{11.} Between the two $\mathrm pK_a$ of the groups that change: $(3.9 + 9.9)/2 \approx 6.9$.
\textbf{12.} In the dipeptide the amine's carbon carries an amide instead of a carboxylate: the
neighbouring groups and charges change, and so do the $\mathrm pK_a$ (the N-terminal ammonium of a
\omterm{def:b2:biomolecules:peptide}{peptide} is markedly more acidic than that of a free amino acid). The model gives the shape, not the
exact values.
\textbf{13.} Acid and amine first form a salt; heating would give mixtures (Asp--Asp, longer chains,
the wrong carboxylic group) and racemise the stereocentres.
\textbf{14.} Its amino group and its side-chain carboxylic group.
\textbf{15.} It turns the \ce{OH} of the acid into a good leaving group, an activated derivative that the
amine attacks at room temperature (acyl substitution).
\textbf{16.} The $\beta$-dipeptide, linked through the \omterm{def:b2:biomolecules:amino-acid}{side chain}: an isomer of aspartame, not aspartame.
\textbf{17.} The amine as a benzyloxycarbonyl carbamate and the \omterm{def:b2:biomolecules:amino-acid}{side chain} as a benzyl ester: both are
removed together by hydrogenolysis over palladium, which leaves the methyl ester untouched (orthogonal
to it).
\textbf{18.} An activated acid has an acidic $\alpha$ hydrogen; base removes it and the stereocentre
racemises.
\textbf{19.} Protect aspartic acid (carbamate on N, benzyl ester on the \omterm{def:b2:biomolecules:amino-acid}{side chain}); activate the free
$\alpha$-\ce{COOH} with the \omterm{def:b2:biomolecules:peptide}{coupling agent}; add the methyl ester of L-phenylalanine; hydrogenolyse.
\textbf{20.} The methyl ester and the \omterm{def:b2:biomolecules:peptide}{peptide bond}; the ester, an amide being far less reactive.
\textbf{21.} The dipeptide Asp--Phe (free acid) and methanol.
\textbf{22.} The sweet molecule is consumed: its hydrolysis products do not replace it.
\textbf{23.} $14 \times 12.011 + 18 \times 1.008 + 2 \times 14.007 + 5 \times 15.999 =
\qty{294.307}{g/mol}$.
\textbf{24.} $0.180/294.307 = \qty{6.116e-4}{mol}$, \qty{0.612}{mmol}.
\textbf{25.} One methanol per aspartame: $0.6116 \times 32.042 = \boldsymbol{\approx \qty{19.6}{mg}}$ of
methanol.
\end{solution}
